\IfFileExists{stacks-project.cls}{%
\documentclass{stacks-project}
}{%
\documentclass{amsart}
}
\usepackage{amssymb}

% For dealing with references we use the comment environment
\usepackage{verbatim}
\newenvironment{reference}{\comment}{\endcomment}
%\newenvironment{reference}{}{}
\newenvironment{slogan}{\comment}{\endcomment}
\newenvironment{history}{\comment}{\endcomment}

% For commutative diagrams we use Xy-pic
\usepackage[all]{xy}

% We use 2cell for 2-commutative diagrams.
\xyoption{2cell}
\UseAllTwocells

% We use multicol for the list of chapters between chapters
\usepackage{multicol}

% This is generally recommended for better output
\usepackage{lmodern}
\usepackage[T1]{fontenc}
\usepackage{textalpha}

% For cross-file-references
\usepackage{xr-hyper}
\newcommand{\maybeexternaldocument}[2]{%
  \IfFileExists{#2.aux}{%
    \externaldocument[#1]{#2}%
  }{%
    \PackageWarning{stack}{Missing #2.aux; external references with prefix #1 unavailable}%
  }%
}
\let\externaldocumentorig\externaldocument
\renewcommand{\externaldocument}[2][]{%
  \IfFileExists{#2.aux}{%
    \externaldocumentorig[#1]{#2}%
  }{%
    \PackageWarning{stack}{Missing #2.aux; external references with prefix #1 unavailable}%
  }%
}

% Package for hypertext links:
\usepackage{hyperref}

% For any local file, say "hello.tex" you want to link to please
% use \externaldocument[hello-]{hello}
\externaldocument[introduction-]{introduction}
\externaldocument[conventions-]{conventions}
\externaldocument[sets-]{sets}
\externaldocument[categories-]{categories}
\externaldocument[topology-]{topology}
\externaldocument[sheaves-]{sheaves}
\externaldocument[sites-]{sites}
\externaldocument[stacks-]{stacks}
\externaldocument[fields-]{fields}
\externaldocument[algebra-]{algebra}
\externaldocument[brauer-]{brauer}
\externaldocument[homology-]{homology}
\externaldocument[derived-]{derived}
\externaldocument[simplicial-]{simplicial}
\externaldocument[more-algebra-]{more-algebra}
\externaldocument[smoothing-]{smoothing}
\externaldocument[modules-]{modules}
\externaldocument[sites-modules-]{sites-modules}
\externaldocument[injectives-]{injectives}
\externaldocument[cohomology-]{cohomology}
\externaldocument[sites-cohomology-]{sites-cohomology}
\externaldocument[dga-]{dga}
\externaldocument[dpa-]{dpa}
\externaldocument[sdga-]{sdga}
\externaldocument[hypercovering-]{hypercovering}
\externaldocument[schemes-]{schemes}
\externaldocument[constructions-]{constructions}
\externaldocument[properties-]{properties}
\externaldocument[morphisms-]{morphisms}
\externaldocument[coherent-]{coherent}
\externaldocument[divisors-]{divisors}
\externaldocument[limits-]{limits}
\externaldocument[varieties-]{varieties}
\externaldocument[topologies-]{topologies}
\externaldocument[descent-]{descent}
\externaldocument[perfect-]{perfect}
\externaldocument[more-morphisms-]{more-morphisms}
\externaldocument[flat-]{flat}
\externaldocument[groupoids-]{groupoids}
\externaldocument[more-groupoids-]{more-groupoids}
\externaldocument[etale-]{etale}
\externaldocument[chow-]{chow}
\externaldocument[intersection-]{intersection}
\externaldocument[pic-]{pic}
\externaldocument[weil-]{weil}
\externaldocument[adequate-]{adequate}
\externaldocument[dualizing-]{dualizing}
\externaldocument[duality-]{duality}
\externaldocument[discriminant-]{discriminant}
\externaldocument[derham-]{derham}
\externaldocument[local-cohomology-]{local-cohomology}
\externaldocument[algebraization-]{algebraization}
\externaldocument[curves-]{curves}
\externaldocument[resolve-]{resolve}
\externaldocument[models-]{models}
\externaldocument[functors-]{functors}
\externaldocument[equiv-]{equiv}
\externaldocument[pione-]{pione}
\externaldocument[etale-cohomology-]{etale-cohomology}
\externaldocument[proetale-]{proetale}
\externaldocument[relative-cycles-]{relative-cycles}
\externaldocument[more-etale-]{more-etale}
\externaldocument[trace-]{trace}
\externaldocument[crystalline-]{crystalline}
\externaldocument[spaces-]{spaces}
\externaldocument[spaces-properties-]{spaces-properties}
\externaldocument[spaces-morphisms-]{spaces-morphisms}
\externaldocument[decent-spaces-]{decent-spaces}
\externaldocument[spaces-cohomology-]{spaces-cohomology}
\externaldocument[spaces-limits-]{spaces-limits}
\externaldocument[spaces-divisors-]{spaces-divisors}
\externaldocument[spaces-over-fields-]{spaces-over-fields}
\externaldocument[spaces-topologies-]{spaces-topologies}
\externaldocument[spaces-descent-]{spaces-descent}
\externaldocument[spaces-perfect-]{spaces-perfect}
\externaldocument[spaces-more-morphisms-]{spaces-more-morphisms}
\externaldocument[spaces-flat-]{spaces-flat}
\externaldocument[spaces-groupoids-]{spaces-groupoids}
\externaldocument[spaces-more-groupoids-]{spaces-more-groupoids}
\externaldocument[bootstrap-]{bootstrap}
\externaldocument[spaces-pushouts-]{spaces-pushouts}
\externaldocument[spaces-chow-]{spaces-chow}
\externaldocument[groupoids-quotients-]{groupoids-quotients}
\externaldocument[spaces-more-cohomology-]{spaces-more-cohomology}
\externaldocument[spaces-simplicial-]{spaces-simplicial}
\externaldocument[spaces-duality-]{spaces-duality}
\externaldocument[formal-spaces-]{formal-spaces}
\externaldocument[restricted-]{restricted}
\externaldocument[spaces-resolve-]{spaces-resolve}
\externaldocument[formal-defos-]{formal-defos}
\externaldocument[defos-]{defos}
\externaldocument[cotangent-]{cotangent}
\externaldocument[examples-defos-]{examples-defos}
\externaldocument[algebraic-]{algebraic}
\externaldocument[examples-stacks-]{examples-stacks}
\externaldocument[stacks-sheaves-]{stacks-sheaves}
\externaldocument[criteria-]{criteria}
\externaldocument[artin-]{artin}
\externaldocument[quot-]{quot}
\externaldocument[stacks-properties-]{stacks-properties}
\externaldocument[stacks-morphisms-]{stacks-morphisms}
\externaldocument[stacks-limits-]{stacks-limits}
\externaldocument[stacks-cohomology-]{stacks-cohomology}
\externaldocument[stacks-perfect-]{stacks-perfect}
\externaldocument[stacks-introduction-]{stacks-introduction}
\externaldocument[stacks-more-morphisms-]{stacks-more-morphisms}
\externaldocument[stacks-geometry-]{stacks-geometry}
\externaldocument[moduli-]{moduli}
\externaldocument[moduli-curves-]{moduli-curves}
\externaldocument[examples-]{examples}
\externaldocument[exercises-]{exercises}
\externaldocument[guide-]{guide}
\externaldocument[desirables-]{desirables}
\externaldocument[coding-]{coding}
\externaldocument[obsolete-]{obsolete}
\externaldocument[fdl-]{fdl}
\externaldocument[index-]{index}

% Heap Project documents
\externaldocument[linear-algebra-]{linear-algebra}
\externaldocument[groups-]{groups}
\externaldocument[complex-analysis-]{complex-analysis}
\externaldocument[functional-analysis-]{functional-analysis}
\externaldocument[categories-]{categories}
\externaldocument[commutative-algebra-]{commutative-algebra}
\externaldocument[ringed-spaces-]{ringed-spaces}
\externaldocument[homological-algebra-]{homological-algebra}
\externaldocument[lie-algebras-]{lie-algebras}
\externaldocument[representation-theory-]{representation-theory}
\externaldocument[smooth-manifolds-]{smooth-manifolds}
\externaldocument[differential-topology-]{differential-topology}
\externaldocument[algebraic-topology-]{algebraic-topology}
\externaldocument[differential-geometry-]{differential-geometry}
\externaldocument[algebraic-geometry-]{algebraic-geometry}
\externaldocument[symplectic-geometry-]{symplectic-geometry}
\externaldocument[complex-geometry-]{complex-geometry}
\externaldocument[hodge-theory-]{hodge-theory}
\maybeexternaldocument{deformations-}{deformations}
\externaldocument[vertex-algebras-]{vertex-algebras}
\externaldocument[springer-]{springer}
\externaldocument[infty-categories-]{infty-categories}
\externaldocument[seminars-differential-geometry-]{%
seminars-differential-geometry}
\externaldocument[seminars-complex-geometry-]{%
seminars-complex-geometry}
\externaldocument[seminars-algebraic-geometry-]{%
seminars-algebraic-geometry}
\externaldocument[seminars-others-]{%
seminars-others}
%\externaldocument[geometric-prequantization-]{geometric-prequantization}
\externaldocument[surfaces-]{surfaces}
\externaldocument[birational-maps-commutative-algebra-]{birational-maps-commutative-algebra}
\externaldocument[cremona-transformations-]{cremona-transformations}
\externaldocument[derived-categories-of-sheaves-]{derived-categories-of-sheaves}
\externaldocument[geometric-invariant-theory-]{geometric-invariant-theory}
\externaldocument[geometric-stability-conditions-and-group-actions-]{geometric-stability-conditions-and-group-actions}
\externaldocument[moduli-spaces-of-sheaves-]{moduli-spaces-of-sheaves}
\externaldocument[bridgeland-stability-general-theory-]{bridgeland-stability-general-theory}
\externaldocument[hilbert-schemes-]{hilbert-schemes}
\externaldocument[noncommutative-abelian-surfaces-and-kummer-type-hyperkahler-varieties-]{noncommutative-abelian-surfaces-and-kummer-type-hyperkahler-varieties}
\externaldocument[mumford-tate-groups-in-hodge-theory-]{mumford-tate-groups-in-hodge-theory}
\externaldocument[hodge-birational-atoms-2026-]{hodge-birational-atoms-2026}
\externaldocument[xxii-egd-2026-]{%
xxii-egd-2026}
\externaldocument[pde-]{pde}
\externaldocument[probability-]{probability}

\externaldocument[linguistics-]{linguistics}
\externaldocument[genealogia-familiar-]{%
genealogia-familiar}

\externaldocument[%
16th-alga-meeting-2026-]{%
16th-alga-meeting-2026}

\externaldocument[%
iv-jornadas-lefschetz-2026-]{%
iv-jornadas-lefschetz-2026}

\externaldocument[reading-the-masters-]{%
reading-the-masters}

\externaldocument[real-analysis-]{%
real-analysis}

\externaldocument[optimal-transport-]{%
optimal-transport}

\externaldocument[history-]{history}

% Theorem environments.
%
\theoremstyle{plain}
\newtheorem{theorem}[subsection]{Theorem}
\newtheorem{proposition}[subsection]{Proposition}
\newtheorem{lemma}[subsection]{Lemma}

\theoremstyle{definition}
\newtheorem{definition}[subsection]{Definition}
\newtheorem{example}[subsection]{Example}
\newtheorem{exercise}[subsection]{Exercise}
\newtheorem{situation}[subsection]{Situation}

\theoremstyle{remark}
\newtheorem{remark}[subsection]{Remark}
\newtheorem{remarks}[subsection]{Remarks}

\numberwithin{equation}{subsection}

% Macros
%
\def\lim{\mathop{\mathrm{lim}}\nolimits}
\def\colim{\mathop{\mathrm{colim}}\nolimits}
\def\Spec{\mathop{\mathrm{Spec}}}
\def\Hom{\mathop{\mathrm{Hom}}\nolimits}
\def\Ext{\mathop{\mathrm{Ext}}\nolimits}
\def\SheafHom{\mathop{\mathcal{H}\!\mathit{om}}\nolimits}
\def\SheafExt{\mathop{\mathcal{E}\!\mathit{xt}}\nolimits}
\def\Sch{\mathit{Sch}}
\def\Mor{\mathop{\mathrm{Mor}}\nolimits}
\def\Ob{\mathop{\mathrm{Ob}}\nolimits}
\def\Sh{\mathop{\mathit{Sh}}\nolimits}
\def\NL{\mathop{N\!L}\nolimits}
\def\CH{\mathop{\mathrm{CH}}\nolimits}
\def\proetale{{pro\text{-}\acute{e}tale}}
\def\etale{{\acute{e}tale}}
\def\QCoh{\mathit{QCoh}}
\def\Ker{\mathop{\mathrm{Ker}}}
\def\Im{\mathop{\mathrm{Im}}}
\def\Coker{\mathop{\mathrm{Coker}}}
\def\Coim{\mathop{\mathrm{Coim}}}

% Boxtimes
%
\DeclareMathSymbol{\boxtimes}{\mathbin}{AMSa}{"02}

%
% Macros for moduli stacks/spaces
%
\def\QCohstack{\mathcal{QC}\!\mathit{oh}}
\def\Cohstack{\mathcal{C}\!\mathit{oh}}
\def\Spacesstack{\mathcal{S}\!\mathit{paces}}
\def\Quotfunctor{\mathrm{Quot}}
\def\Hilbfunctor{\mathrm{Hilb}}
\def\Curvesstack{\mathcal{C}\!\mathit{urves}}
\def\Polarizedstack{\mathcal{P}\!\mathit{olarized}}
\def\Complexesstack{\mathcal{C}\!\mathit{omplexes}}
% \Pic is the operator that assigns to X its picard group, usage \Pic(X)
% \Picardstack_{X/B} denotes the Picard stack of X over B
% \Picardfunctor_{X/B} denotes the Picard functor of X over B
\def\Pic{\mathop{\mathrm{Pic}}\nolimits}
\def\Picardstack{\mathcal{P}\!\mathit{ic}}
\def\Picardfunctor{\mathrm{Pic}}
\def\Deformationcategory{\mathcal{D}\!\mathit{ef}}


\begin{document}

\title{African languages in Brazil}
\maketitle

\phantomsection
\label{section-phantom}

\tableofcontents

\section{Method}
\label{section-african-languages-method}

This notebook began with words encountered
at MUNCAB and the Casa do Benin on
September 26, 2026.

The first methodological lesson is that
``African word'' is not a sufficient
classification.  The visit produced words
from Yoruba-related traditions and from
Bantu-language histories, as well as
Brazilian diasporic labels.

Three forms should be distinguished:

\begin{itemize}
\item the spelling printed by the museum;
\item a Brazilian religious spelling;
\item a fully tone-marked spelling in the
African language.
\end{itemize}

The last form requires underdots and tone
marks.  The present pdfLaTeX setup does not
yet represent every combined Yoruba
diacritic elegantly.  Museum spellings are
therefore preserved first, with linguistic
normalization stated separately and
cautiously.

\section{Yoruba and Nago}
\label{section-languages-yoruba-nago}

Yoruba names a language and a broad
historical and cultural field in West
Africa, especially in present-day Nigeria
and Benin.

Nag\^o is a diasporic category used in
Brazil, especially Bahia, for people and
traditions strongly connected with the
Yoruba world.  It is not simply the modern
African self-name Yoruba.

The word is usually related to forms such
as \emph{Anago}, used in the Bight of Benin
for Yoruba-speaking populations.  The
historical extension of the term changed
through the Atlantic trade and in Brazil
\cite{lovejoyojo2015lucumi}.

In Candombl\'e, linguistic practice is one
important marker of the ritual categories
called nations \cite{castro1981lingua}.
The museum map therefore said Yoruba in
Africa, while the Afonj\'a book spoke of a
Nag\^o cosmology in Bahia.

\section{Omi Tutu}
\label{section-languages-omi-tutu}

The MUNCAB exhibition title was printed as
\emph{Omi Tutu---\'Agua fresca}.

The first word, \emph{omi}, means water.
The form \emph{tutu} belongs to the field
of coolness or freshness.  In religious
and aesthetic contexts, coolness can also
carry ideas of calm, balance, and
appeasement.

The title should therefore not be reduced
to the temperature of water.  The museum
itself connected fresh water with vitality,
creation, conciliation, and abundance
\cite{muncab2026omitutu}.

\section{Ori Eja and Iya Eja}
\label{section-languages-ori-iya-eja}

Two labels beside Hariel Revignet's work
were printed as:

\begin{itemize}
\item \emph{Ori Ej\'a};
\item \emph{Iya Ej\'a}.
\end{itemize}

A first lexical analysis gives:

\begin{itemize}
\item \emph{ori}: head, with a wider field
that can include personal destiny or the
inner spiritual head;
\item \emph{eja}: fish;
\item \emph{iya}: mother.
\end{itemize}

This suggests readings such as ``fish
head'' or ``head of fish'', and ``mother of
fish''.  These are lexical decompositions,
not authoritative translations of the
artist's titles.  Yoruba tone and syntax
must be checked before fixing the full
phrases.

The word \emph{ori} is especially
important.  In the religious field it
cannot be reduced to anatomy.  It can name
a dimension of personhood and destiny.

\section{Orun, Aiye, and Olorun}
\label{section-languages-orun-aiye}

The film title
\emph{Orun Aiye---A cria\c{c}\~ao do mundo}
made three similar-looking forms easy to
confuse.

\emph{Orun} names the spiritual or
invisible realm.  \emph{Aye} or
\emph{Aiye} names the visible earthly
world.  The two realms are related rather
than simply sealed off from one another.

\emph{Olorun} is a name or title for the
supreme divinity.  It is not another
spelling of the realm \emph{Orun}.

This distinction explains the museum
phrase that M\~ae Stella, after death,
returned to Orun.  It does not say that she
became Olorum.

\section{Ile and Ile Aiye}
\label{section-languages-ile-aiye}

\emph{Ile} can mean house or home.  It
appears in names of religious houses, such
as Il\^e Ax\'e Op\^o Afonj\'a.

It also appears in the name Il\^e Aiy\^e,
the bloco afro.  The exact semantics of a
proper name should not be produced by
mechanically joining dictionary glosses,
but the words place house, earth, and world
inside its semantic field.

\section{Ofa}
\label{section-languages-ofa}

The iron sculpture encountered beside the
Diogum biography represented an
\emph{ofa}, the bow and arrow associated
with Ox\'ossi.

The standard Yoruba form is commonly
written with an underdotted first vowel and
a low tone on the final vowel.  The
Brazilian form \emph{of\'a} records the word
without the full Yoruba orthography.

This is a good example of a word that moved
with religious knowledge and was adapted
to Brazilian Portuguese spelling.

\section{Iyami}
\label{section-languages-iyami}

The N\'adia Taquary panel printed
\emph{Iy\'amis} and translated the term as
``feiticeiras''.

The beginning of the word is related to
the Yoruba word for mother.  In the
religious concept, however, ancestral
maternal and feminine power cannot be
exhausted by the English ``witch'' or the
Portuguese ``feiticeira''
\cite{santos2008iyami}.

The encounter suggests a translation
problem, not only a vocabulary item.  What
features of a concept disappear when a
European category of witchcraft is used?

The comparison with Circe should remain a
comparative literary question.  Similar
images of feared or transformative female
power do not establish common origin.

\section{Funfun Dudu}
\label{section-languages-funfun-dudu}

The title of the Alberto Pitta book was
printed as \emph{F\'unf\'un D\'ud\'u}.

The two Yoruba colour terms are associated
with white and black or dark.  The title is
especially suggestive for an artist whose
work moves among textile, print, painting,
carnival, terreiro, and city
\cite{barata2025pitta}.

The linguistic note should not replace the
art-historical one.  The meaning of the
book title belongs to the visual project,
not only to two dictionary entries.

\section{Mironga}
\label{section-languages-mironga}

The Casa do Benin exhibition used
\emph{mironga} in the fields of mystery,
secret, guarded knowledge, spell, and
strategy.  It also used the verb
\emph{mirongar} for working an image through
what can be shown and what remains hidden.

Aulete treats \emph{mironga} as a variant
of \emph{milonga} and relates the latter to
a Kimbundu form, \emph{mi'lona}, glossed as
``word'' \cite{aulete-milonga}.

This is evidence for a Bantu or Angolan
linguistic layer in Brazilian Portuguese,
not a Yoruba one.

The exact historical morphology and the
semantic path from ``word'' to mystery,
spell, or strategic concealment still need
a specialist Kimbundu source.  The note
must preserve that uncertainty.

\section{Zumbi and Nzambi}
\label{section-languages-zumbi-nzambi}

The statue of Zumbi dos Palmares prompted
a memory of the name Nzambi or Nzambi a
Mpungu.

Zumbi is the name by which a historical
leader of Palmares is known.  Nzambi is a
divine name in Kongo and other Central
African religious contexts.

The two words should not be identified
only because they sound similar.  Their
possible etymological histories require
specialist Bantu historical linguistics.

\section{Angola as two kinds of name}
\label{section-languages-angola}

Angola is the name of a modern African
state.  Angola is also the name of a
Candombl\'e nation in Brazil.

The second use emerged through histories
of Central African peoples and Bantu
languages, including Kimbundu and Kongo
fields.  It does not mean that the ritual
nation is simply the modern country
transplanted intact.

This distinction is one instance of a
larger rule: modern political borders,
ethnolinguistic regions, and Brazilian
diasporic categories are different maps.

\section{Akan and West Africa}
\label{section-languages-akan}

A MUNCAB mask label said ``Povos Akan---
\'Africa Central''.  A map in the same
museum placed the Akan region in Ghana and
C\^ote d'Ivoire.

The latter location is West Africa.  The
internal contradiction should be recorded
rather than silently repaired.

Akan is not another word for Yoruba, Jeje,
or a Candombl\'e nation.  It names a large
West African ethnolinguistic field that
includes, among others, Asante and Fante
histories.

\section{Terms still to investigate}
\label{section-languages-open-terms}

The next linguistic work should include:

\begin{itemize}
\item a reliable Yoruba dictionary for all
tones and underdots;
\item the syntax of the titles
\emph{Ori Eja} and \emph{Iya Eja};
\item the relation among Yoruba, Anago,
Nag\^o, and Cuban Lucum\'\i;
\item a specialist Kimbundu etymology of
\emph{mironga} or \emph{milonga};
\item the history of Nzambi and the name
Zumbi;
\item African-language layers in Brazilian
Portuguese beyond Yoruba and Kimbundu.
\end{itemize}

\bibliography{my,salvador}
\bibliographystyle{amsalpha}

\end{document}
